% !TeX root = ../main.tex

\chapter{Evaluation}

\section{Evaluation Goals}

The evaluation is designed to answer whether the workflow produces reports
that are useful for auditors. This differs from a pure detector evaluation. A
detector evaluation asks whether the system predicts that a project is
vulnerable. An audit-usefulness evaluation asks whether the system points to
the right code, explains the right root cause, and avoids excessive false
reports.

The evaluation therefore has three goals:

\begin{enumerate}
  \item Measure protocol-level detection performance for compatibility with
        prior work.
  \item Measure report-level correctness, including localization and
        explanation quality.
  \item Measure over-reporting, because unsupported extra findings increase
        manual audit cost.
\end{enumerate}

\section{Dataset}

The evaluation dataset is built from public audit findings. Each record
contains
metadata such as report date, audit firm, severity, protocol, title, tags,
year, source-code link, Solodit URL, GitHub link, summary, whether the issue is
classified as price, reward, or vote manipulation, whether source code is
available, and remarks.

The dataset is organized into three subsets.

\textbf{Collected findings} include all candidate findings from the collection
pipeline. This set is useful for understanding the distribution of protocols,
firms, severities, and issue categories, but it may contain entries that are
not directly analyzable.

\textbf{Source-available findings} include records with public source code.
This set is closer to the target evaluation but still may include repositories
that are difficult to build or incomplete.

\textbf{Build-verified findings} include records whose repositories can be
compiled or analyzed reproducibly. This is the primary subset for quantitative
evaluation because the static-analysis front end depends on buildable Solidity
projects.

This separation avoids overstating the dataset size. A large collected dataset
does not automatically imply a large reproducible evaluation set.

\section{Metrics}

\subsection{Why Binary Metrics Are Insufficient}

Precision, recall, and F1-score are standard metrics for vulnerability
detection. They are easy to compute and useful for comparing classifiers. For
audit assistance, however, they are incomplete. A binary true positive can hide
multiple report-level failures that matter in practice.

First, a system may correctly flag a project as vulnerable but still point to
the wrong code path. For an auditor, this is much less useful than a report
that identifies the relevant source, sink, or affected function. Second, a
system may emit too many vulnerable groups or findings. If only one report is
correct, binary scoring may still reward the system while ignoring the extra
manual review cost of unsupported reports. Third, a system may localize roughly
the right region but explain the wrong root cause or recommend the wrong
mitigation. In that case, the output still requires substantial human
reinterpretation before it becomes audit-useful.

These limitations motivate report-level evaluation alongside protocol-level
metrics. In this thesis, a generated report is therefore additionally judged by
localization accuracy, explanation quality, mitigation agreement, semantic
similarity to the corresponding audit finding, and the number of unsupported
extra reports. Binary metrics remain useful as a baseline, but they are not
treated as the final measure of success.

\subsection{Binary Metrics}

The first metric group follows standard vulnerability detection evaluation. A
project-level prediction is positive if the system emits at least one
price-manipulation finding for the project. Ground truth is positive if the
project contains at least one relevant audit finding. From these labels, the
evaluation computes true positives, false positives, false negatives,
precision, recall, and F1-score.

These metrics are reported mainly for compatibility with prior work such as
PMDetector \citep{pmdetector}. Their role in this thesis is to provide a
baseline view of detector behavior before moving to more detailed report-level
measurements.

\subsection{Localization Metrics}

The second metric group evaluates whether a generated report points to the
right place. Localization can be measured at several levels:

\begin{itemize}
  \item Protocol-level localization: the prediction belongs to the same
        project as the ground-truth finding.
  \item Contract or function localization: the predicted root-cause location
        names the same contract, function, wrapper, or semantically equivalent
        code region.
  \item Source localization: the predicted vulnerable read corresponds to the
        same price source or source wrapper.
  \item Sink localization: the predicted sink corresponds to the same economic
        operation as the audit finding.
  \item Mitigation localization: the predicted missing mitigation matches the
        ground-truth mitigation category.
\end{itemize}

These levels make partial success visible. For example, a prediction may
identify the correct spot price source but miss the exact liquidation sink.
Another prediction may identify the affected function but give the wrong
mitigation.

\subsection{Semantic Explanation Score}

The explanation score compares the generated explanation with the audit-finding
summary. The evaluation embeds both texts using a sentence embedding model and
computes cosine similarity:

\[
\mathrm{DescriptionScore}(p, g) =
\begin{cases}
\cos(E(p), E(g)), & \cos(E(p), E(g)) > \tau \\
0, & \text{otherwise}
\end{cases}
\]

where $p$ is the predicted explanation, $g$ is the ground-truth explanation,
$E(\cdot)$ is the embedding function, and $\tau$ is the similarity threshold.
In this thesis, the threshold is set to 0.7, inspired by semantic evaluation
settings such as the Bastet competition \citep{bastet}. This fixed threshold
provides a consistent operating point for report matching, although its
sensitivity remains a validity consideration discussed later in this chapter.

\subsection{Over-Reporting Penalty}

The over-reporting penalty measures how many extra reports the system emits for
a project:

\[
\mathrm{Penalty}(p) =
\max(0, \#\mathrm{Predictions}(p) - \#\mathrm{GroundTruth}(p)).
\]

This penalty is important because the static-analysis phase intentionally
over-approximates candidate paths. Without a penalty, the system could improve
recall by emitting many weak findings. In an audit setting, such behavior is
costly because each report must be checked manually.

\section{Matching Procedure}

For each project, the evaluation builds a pairwise score matrix between
predicted reports and ground-truth audit findings. Each pair score combines
localization agreement, mitigation agreement, and semantic explanation
similarity. Scores below a minimum threshold are set to zero. Maximum-weight
bipartite matching is then used to assign predictions to ground-truth findings.

After matching, each prediction is classified as matched or unmatched, and each
ground-truth finding is classified as covered or missed. The evaluation reports
the total matched score, unmatched predictions, missed findings, binary metrics,
and over-reporting penalty.

\section{Baselines}

The main baseline is a PMDetector-style binary evaluation of the current
workflow. Under this baseline, a project is counted as correctly detected if at
least one relevant vulnerability is found. This baseline intentionally ignores
whether the report localizes the right path or explains the right root cause.

The evaluation framework also includes a simpler PMDetector-style prompt
baseline that emits one report containing
\texttt{is\_high\_risk}, vulnerable functions, vulnerable groups, and a reason.
Comparing this baseline with the staged workflow would show whether decomposing
source classification, sink triage, and root-cause judgment improves
audit-useful output.

The published PMDetector results should be treated as prior-work context rather
than a direct baseline unless the dataset and implementation are aligned. The
current thesis uses different ground truth and different evaluation metrics.

\section{Qualitative Results}

The implemented workflow produces analyzable artifacts for several projects.
This section presents representative qualitative results that illustrate both
the strengths and the limitations of the workflow.

\subsection{Spartan}

For the Spartan project, one analysis summary reports 25 analyzed contracts, 67
raw source candidates, 17 canonical sources, 33 sink sites, 217 raw vulnerable
paths, and 97 representative paths. The LLM report filters eight source labels
down to two retained spot-price sources: \texttt{iPOOL.baseAmount} and
\texttt{iPOOL.tokenAmount}. It then marks 13 sink groups as dangerous and emits
two vulnerable root-cause findings.

This result illustrates the intended pipeline behavior. Static analysis
generates many raw paths, representative grouping reduces redundancy, and the
LLM collapses evidence into a smaller number of root-cause findings. The
key evaluation question is whether those findings align with the corresponding
audit or incident ground truth.

\subsection{Predy}

For Predy, the LLM report begins from 41 representative paths and six source
labels. It keeps \texttt{IUniswapV3PoolState.slot0} and
\texttt{Trade.getSqrtPrice}, drops fee, liquidity, TWAP, and protocol-state
sources, and emits two vulnerable root-cause findings. The generated
explanations point to spot-price use in \texttt{Perp.reallocate} and
\texttt{Trade.trade}, with missing TWAP or slippage checks.

This result shows the value of source classification. Without the first LLM
stage, liquidity and protocol-state values could become false positives. The
stage decomposition helps focus root-cause judgment on price-like sources.

\subsection{Abracadabra Money}

For one Abracadabra Money run, the summary reports 87 analyzed contracts, 56
raw source candidates, and 23 canonical sources, but zero sink sites and zero
representative paths. This is a useful failure case. It may indicate that the
current sink definitions do not cover the relevant economic operation, that the
selected analysis target does not include the affected code path, or that the
issue requires modeling that is not yet implemented.

This case supports the need for transparent reporting. A binary evaluation
would count the project as a miss if ground truth contains a relevant finding,
but the engineering diagnosis is more specific: source identification may work,
while sink coverage fails.

\section{Threats to Validity}

\textbf{Construct validity.} Semantic similarity may not perfectly represent
auditor judgment. The embedding model and threshold therefore approximate, but
do not fully replace, human assessment.

\textbf{Internal validity.} LLM results may depend on model version, prompt
wording, temperature, and chunking. These dependencies should be considered
when interpreting the reported results.

\textbf{External validity.} The thesis focuses on price manipulation in
Solidity DeFi protocols. Results may not generalize to other vulnerability
classes, non-EVM chains, or protocols with off-chain components.

\textbf{Dataset validity.} Audit findings differ in specificity and quality.
Some findings are disputed, duplicated, or only partially related to price
manipulation. The dataset therefore requires explicit inclusion and exclusion
criteria.

\textbf{Implementation validity.} The current implementation may miss paths because of
dynamic dispatch, proxy patterns, incomplete build targets, missing sink
definitions, or unsupported Solidity patterns. These failures are reported
explicitly rather than hidden.

\section{Evaluation Interpretation}

The main contribution of this evaluation design is a clearer measurement of
audit usefulness. A system may have good protocol-level detection while still
showing poor localization, weak explanation matching, or high over-reporting;
in such cases, binary metrics alone are misleading. Conversely, stronger
report-level matching provides direct evidence that staged source and sink
reasoning improves audit-useful output. This interpretation supports the
central thesis that LLM-assisted security tools should be evaluated at the
report level rather than only through binary labels.
